\documentclass{article}
\usepackage[utf8]{inputenc}
\usepackage{amsmath}
\usepackage{amssymb}
\usepackage{algorithm}
\usepackage{algpseudocode}
\usepackage{graphicx}
\usepackage{geometry}
\geometry{a4paper, margin=1in}




\newcommand{\old}[1]{{}}
\newtheorem{theorem}{Theorem}
\newtheorem{corollary}{Corollary}
\newtheorem{conjecture}[theorem]{Conjecture}
\newtheorem{lemma}[theorem]{Lemma}
\newtheorem{claim}[theorem]{Claim}
\newtheorem{observation}[theorem]{Observation}
\newtheorem{definition}[theorem]{Definition}

\title{Descending Greedy Filter}
\author{P. Carmi, N. Horovitz, Y. Yashuk}
\date{March 2026}




\begin{document}

\maketitle

\begin{abstract}
 We present a new construction of geometric spanners that demonstrates improved empirical performance over the classical greedy spanner on random point sets. For a range of stretch parameters $t>1$, our method consistently produces spanners with fewer edges and lower total weight while maintaining the desired stretch. To the best of our knowledge, this is the first construction that empirically outperforms the greedy spanner on random point sets in terms of both sparsity and total weight. Although the greedy spanner is widely regarded as the most effective practical construction due to its strong theoretical properties, our experimental results indicate that, on typical random inputs, the proposed construction yields significantly sparser and lighter spanners.
\end{abstract}



\section{Introduction}
Geometric spanners are a central object in computational geometry. Given a set $P$ of points in the plane and a stretch parameter $t \ge 1$, a geometric $t$-spanner is a sparse graph in which the shortest-path distance between any pair of points is at most $t$ times their Euclidean distance. Such structures have been extensively studied due to their applications in network design, approximation algorithms, and distributed systems.

Among the many known constructions, the greedy spanner is widely regarded as the most effective in practice, achieving a favorable trade-off between sparsity, total weight, and stretch. 

In this work, we present a new construction of geometric spanners and evaluate its performance on random point sets. For a range of stretch parameters $t>1$, our method consistently produces spanners with fewer edges and lower total weight than the greedy spanner, while maintaining the desired stretch. To the best of our knowledge, this is the first construction that empirically outperforms the greedy spanner on random point sets in terms of both sparsity and total weight.

In addition to its empirical advantages, our construction shares a key structural property with the greedy spanner: it is edge-minimal with respect to the $t$-spanner condition. That is, removing any edge from the resulting graph violates the stretch requirement. As we discuss later, this property is not satisfied by most classical constructions, such as $\Theta$-graphs, Yao graphs, and WSPD-based spanners.

Our results suggest that, for typical input distributions, it is possible to obtain spanners that are both sparser and lighter than those produced by the greedy algorithm, while preserving strong structural guarantees.


\section{The algorithm}

\algrenewcommand{\algorithmiccomment}[1]{\hspace{3em}// #1}
\begin{algorithm}
\caption{Descending Greedy Filter}
\begin{algorithmic}[1]
\Statex \textbf{Input:} 
\Statex \hspace{1cm} $P$: a set of $n$ points in $\mathbb{R}^d$.
\Statex \hspace{1cm} $t$: a real constant s.t. $t > 1$.
%\Statex \hspace{1cm} $E$: a set of edges to filter, default to the complete graph over $P$.
\Statex \textbf{Output:} 
\Statex \hspace{1cm} a set of edges, $E$ s.t. $(P,E)$ is a $t-spanner$ for $P$.

\State sort the $n \choose {2} $ pairs of distinct points in \textbf{non-ascending} order of their distances

(breaking ties arbitrarily) and store them in a list $L$.
\State $E \leftarrow L$ 
\For{$(p,q) \in L$ in \textbf{non-ascending} order}
    \State $E \leftarrow E \setminus (p,q)$
    \For{$(r,s) \in P \times P$}
        \If{$\delta_{E}(s,r) > t \cdot |sr|$}
            \State $E \leftarrow E \cup (p,q)$
            \State \textbf{break}
        \EndIf
    \EndFor
\EndFor
\State \Return $E$
\end{algorithmic}
\end{algorithm}


\clearpage


\section{Structural Property}
A key structural property shared by our construction and the greedy spanner is edge-minimality with respect to the $t$-spanner condition. In both constructions, every edge is essential: removing any edge results in a graph that no longer satisfies the $t$-spanner property. In the greedy spanner, this follows directly from its construction, as each edge is added only when no existing path of stretch at most $t$ connects its endpoints, and thus remains necessary in the final graph. Our construction exhibits the same property—despite being derived from a different principle, every edge is required to preserve the stretch guarantee, and its removal necessarily violates the $t$-spanner condition. To the best of our knowledge, no other classical geometric spanner construction guarantees this edge-minimality property. This highlights a strong structural similarity between the two constructions, despite their significant differences in design and empirical performance.

Most classical geometric spanner constructions do not guarantee edge-minimality. In particular, well-known approaches such as $\Theta$-graphs, Yao graphs, and WSPD-based spanners construct edges according to predetermined geometric or combinatorial rules rather than necessity with respect to the $t$-spanner condition. As a result, these constructions typically include redundant edges, in the sense that some edges can be removed without violating the stretch requirement.

A common optimization in the construction of geometric spanners is to first compute a $(t/t')$-spanner (for example, using a $\Theta$-graph, Yao graph, or WSPD-based construction), and then apply the greedy algorithm with parameter $t'$ on this sparse subgraph. While this approach reduces the number of candidate edges, it does not, in general, preserve the edge-minimality property. Indeed, since the greedy procedure is restricted to a predefined subset of edges, it may include edges that would be redundant if all pairwise edges were available, and thus the resulting graph need not be edge-minimal.

For example, let $a=(0,\frac{1}{2})$, $b=(0,- \frac{1}{2})$ $c=(\frac{\sqrt{t^2-1}}{2} ,0)$, and $d=(\frac{\sqrt{(3t+1)(5t-1)}}{2(t-1)} ,0)$.


\section{Filtering a $t-spanner$}

\begin{algorithm}
\caption{Descending Greedy Filter on t-spanner}
\begin{algorithmic}[1]
\Statex \textbf{Input:} 
\Statex \hspace{1cm} $P$: a set of $n$ points in $\mathbb{R}^d$.
\Statex \hspace{1cm} $t$: a real constant s.t. $t > 1$.
\Statex \hspace{1cm} $E_i$: a $t-spanner$ for P.
\Statex \textbf{Output:} 
\Statex \hspace{1cm} a set of edges, $E$ s.t. $E \subseteq E_i$, $(P,E)$ is a $t-spanner$ for $P$.

\State sort $E_i$ in \textbf{non-ascending} order of their distances

(breaking ties arbitrarily) and store them in a list $L$.
\State $E \leftarrow L$ 
\For{$(p,q) \in L$ in \textbf{non-ascending} order}
    \State $E \leftarrow E \setminus (p,q)$
    \For{$(r,s) \in P \times P$}
        \If{$\delta_{E}(s,r) > t \cdot |sr|$}
            \State $E \leftarrow E \cup (p,q)$
            \State \textbf{break}
        \EndIf
    \EndFor
\EndFor
\State \Return $E$
\end{algorithmic}
\end{algorithm}

We will prove that if fed with $E_i$, a $t-spanner$, the output $E$ of DGF will be:
\begin{enumerate}
    \item $E$ is a $t-spanner$ for $P$.
    \item $E$ is $t-minimal$.
\end{enumerate}

\subsection{Proof:}

The first property is trivial.\\
the second:\\
let $(p,q) \in E$, need to prove: $\exists (r, s) \in P^2, \delta_{E \setminus (p,q)}(r,s) > t \cdot |r s|$.\\
denote $E'$ as the state of $E$ when DGF inspected $(p,q)$.\\
observe $E \subseteq E'$, because the algorithm only removes edges.\\
because $(p,q) \in E$ then  $\exists (r, s) \in P^2, \delta_{E' \setminus (p,q)}(r,s) > t \cdot |r s|$.\\
because $E \subseteq E'$, than  $ \delta_{E \setminus (p,q)}(r,s) \geq \delta_{E' \setminus (p,q)}$.\\
and we get:\\
$\delta_{E \setminus (p,q)}(r,s) \geq \delta_{E' \setminus (p,q)}(r,s) > t \cdot |r s|$\\
so $\delta_{E \setminus (p,q)}(r,s) > t \cdot |r s|$, as needed.\\

\section{Proofs}


\subsection{Bounded Weight}

Denote $G_{dgf} = (P, E_f)$ as the output of the Descending Greedy Filter algorithm. \\
Denote $n = |P|$.

\subsubsection{If $t=n-1$ than $\text{MST}(P) = \text{DGF}(P)$}

Assume $t = n-1$, need to prove:
\begin{enumerate}
    \item $\text{MST}(P) \subseteq \text{DGF}(P)$
    \item $\text{DGF}(P) \subseteq \text{MST}(P)$
\end{enumerate}

\textbf{Prove $\text{MST}(P) \subseteq \text{DGF}(P)$:}\\
Let $(a,b) \in \text{MST}(P)$; need to prove $(a,b) \in \text{DGF}(P)$.\\
Assume for contradiction $(a,b) \notin \text{DGF}(P)$, so when DGF checked $(a,b)$ there was a path:\\
$R = (a = v_1, v_2, \ldots, v_k = b)$ such that:

\begin{enumerate}
    \item[(a)] $\forall i \in [1, k-1] : |v_i v_{i+1}| \leq |ab|$
    \item[(b)] $\mathrm{wt}(R) = \sum_{i=1}^{k-1} |v_i v_{i+1}| \leq t \cdot |ab|$
\end{enumerate}
So when using Kruskal's algorithm for finding the MST (needs citation), when checking $(a,b)$, the vertices $a$ and $b$ must already be in the same connected component.
Therefore it will not choose $(a,b)$ for $\text{MST}(P)$, a contradiction.
\\\\

\textbf{Prove $\text{DGF}(P) \subseteq \text{MST}(P)$:}\\
Let $(a,b) \in \text{DGF}(P)$; we need to prove $(a,b) \in \text{MST}(P)$.\\
Because $(a,b) \in \text{DGF}(P)$, there exist $x, y \in P$ such that
\[
    \delta_{E_f \setminus (a,b)}(x,y) > t \cdot |xy|.
\]
Take these $x, y$ and denote $R = \delta_{E_f \setminus (a,b)}(x,y)$.

We will prove that $|ab| \leq |xy|$.
Assume for contradiction $|xy| < |ab|$; then when DGF inspected $(a,b)$, the edge $(x,y)$ was still in $E_f$ and must have been valid.

We divide into two cases:

\textbf{Case (a):} $\mathrm{wt}(R) < \infty$.

Then $R = (x = v_1, v_2, \ldots, v_k = y)$. Since $R$ is a shortest path it is simple, so $k \leq n-1$. Therefore
\[
    \mathrm{wt}(R) = \sum_{i=1}^{k-1} |v_i v_{i+1}| \leq k \cdot |ab| \leq (n-1)|ab| \leq (n-1) \cdot |xy| = t \cdot |xy|.
\]
So $R$ is a valid path between $x$ and $y$ not using $(a,b)$, a contradiction. This case is impossible.

\textbf{Case (b):} $\mathrm{wt}(R) = \infty$.

Then $x$ and $y$ are in different connected components; denote them as $X$ and $Y$.

We will prove that when running Prim's algorithm for finding the MST (needs citation), the edge $(a,b)$ will be selected.

When DGF removed $(a,b)$ to inspect it, $X$ and $Y$ became disconnected. Hence all other edges between $X$ and $Y$ had already been removed by DGF, meaning they are all longer than $(a,b)$. Therefore Prim's algorithm must select $(a,b)$.

So $(a,b) \in \text{MST}(P)$.

\subsection{t-minimal Conjecture}

\subsubsection{t-minimal Property Definition:}

A set of edges $E$ conforms to the $t-minimal$ property if:
\begin{enumerate}
    \item $E$ is a $t-spanner$ for $P$.
    \item for all $(p,q) \in E$, exists a pair $r,s \in P^2$ such that $\delta_{E \setminus (p,q)}(r,s) > t \cdot |x y|$\\
\end{enumerate}
Or in a formal mathematical writing:\\
\[
\forall (p,q) \in E, \exists (r, s) \in P^2, \delta_{E \setminus (p,q)}(r,s) > t \cdot |x y|.
\]

\begin{conjecture}
    If $E$ is $t-minimal$, then $wt(E) = O(wt(MST(P)))$\\
\end{conjecture}
%
\begin{conjecture}
    $wt(DGF(P)) = O(wt(MST(P)))$\\
\end{conjecture}


\section{Faster Implementation via Binary Search}

The na\"ive implementation of the Descending Greedy Filter (Algorithm~1) iterates over all edges one by one, and for each edge performs an all-pairs shortest path (APSP) computation to check whether its removal causes a violation. This results in $\Theta(m)$ APSP calls, where $m = \binom{n}{2}$ is the total number of edges. We observe that, in practice, the vast majority of edges are removable and only a small subset is necessary. This motivates a binary search approach that can certify large groups of removable edges with a single APSP call.

\subsection{Key Observation}

Let $E$ be the current edge set sorted in non-ascending order by weight, and let $S \subseteq E$ be a contiguous batch of edges. If removing all edges of $S$ from $E$ causes no violation of the $t$-spanner condition (i.e., for all $r, s \in P$, $\delta_{E \setminus S}(r,s) \le t \cdot |rs|$), then every edge in $S$ is removable. This allows us to eliminate an entire batch with a single APSP computation. When a violation is detected, we split the batch and recurse on each half to identify the necessary edges.

\subsection{Algorithm}

\begin{algorithm}
\caption{Descending Greedy Filter with Binary Search}
\begin{algorithmic}[1]
\Statex \textbf{Input:}
\Statex \hspace{1cm} $P$: a set of $n$ points in $\mathbb{R}^d$.
\Statex \hspace{1cm} $t$: a real constant s.t.\ $t > 1$.
\Statex \textbf{Output:}
\Statex \hspace{1cm} A set of edges $E$ s.t.\ $(P, E)$ is a $t$-spanner for $P$.

\State Sort all $\binom{n}{2}$ pairs in \textbf{non-ascending} order by distance into a list $L$.
\State $E \leftarrow L$
\State $E \leftarrow \textsc{Bisect}(E, L, P, t)$
\State \Return $E$
\end{algorithmic}
\end{algorithm}

\begin{algorithm}
\caption{\textsc{Bisect}$(E, S, P, t)$}
\begin{algorithmic}[1]
\Statex \textbf{Input:}
\Statex \hspace{1cm} $E$: current edge set.
\Statex \hspace{1cm} $S$: a batch of candidate edges to remove (already removed from $E$), sorted in non-ascending order by weight.
\Statex \hspace{1cm} $P, t$: point set and stretch parameter.
\Statex \textbf{Output:}
\Statex \hspace{1cm} Updated edge set $E$ with removable edges from $S$ discarded.

\If{$S = \emptyset$}
    \State \Return $E$
\EndIf
\State Compute APSP on $(P, E)$
\If{$\forall (r,s) \in P \times P$: $\delta_E(r,s) \le t \cdot |rs|$}
    \State \Return $E$ \Comment{No violation: all edges in $S$ are removable}
\EndIf
\If{$|S| = 1$}
    \State $E \leftarrow E \cup S$ \Comment{Single necessary edge: restore it}
    \State \Return $E$
\EndIf
\State Split $S$ into $S_1$ (first half) and $S_2$ (second half)
\State $E \leftarrow E \cup S$ \Comment{Restore all edges in $S$}
\State $E \leftarrow E \setminus S_1$
\State $E \leftarrow \textsc{Bisect}(E, S_1, P, t)$ \Comment{Recurse on first half (longer edges)}
\State $E \leftarrow E \setminus S_2$
\State $E \leftarrow \textsc{Bisect}(E, S_2, P, t)$ \Comment{Recurse on second half (shorter edges)}
\State \Return $E$
\end{algorithmic}
\end{algorithm}

\subsection{Correctness}

The binary search implementation produces the same output as the na\"ive algorithm. Each edge is either certified as removable (when the entire batch containing it passes the violation check) or identified as necessary (when the recursion reaches a single-edge base case with a violation). Since the edges are processed in the same non-ascending order and the removal/restoration operations are symmetric, the final edge set is identical to that of Algorithm~1.

\subsection{Time Complexity}

Let $m$ denote the number of input edges and let $k$ denote the number of necessary edges in the output (i.e., $k = |E|$ in the final spanner). Let $T_{\mathrm{APSP}}(n)$ denote the cost of a single APSP computation.

\paragraph{APSP calls.}
Each necessary edge induces a root-to-leaf path in the recursion tree of depth at most $\lceil \log_2 m \rceil$, contributing $O(\log m)$ APSP calls. Each such path terminates at a single-edge base case where the edge is restored. Conversely, groups of removable edges are eliminated at internal nodes without further recursion. The total number of APSP calls is therefore
\[
    O(k \log m),
\]
where $k$ is the number of necessary edges.

\paragraph{Overall complexity.}
Sorting the edges takes $O(m \log m)$ time. The binary search performs $O(k \log m)$ APSP calls. Each APSP computation uses Dijkstra from all $n$ sources, costing $O(n(m + n \log n))$. In addition, each level of the recursion performs $O(m)$ total work for edge removal and restoration in the CSR data structure. The total time complexity is therefore
\[
    O\!\left(m \log m + k \log m \cdot n(m + n \log n)\right).
\]

\paragraph{Comparison with the na\"ive approach.}
The na\"ive implementation performs exactly $m$ APSP calls, for a total cost of $O(m \cdot n(m + n \log n))$. The binary search reduces this to $O(k \log m \cdot n(m + n \log n))$. Since geometric $t$-spanners have $k = O(n)$ edges while $m = \Theta(n^2)$, the speedup factor is
\[
    \frac{m}{k \log m} = \frac{\Theta(n^2)}{O(n \log n)} = \Omega\!\left(\frac{n}{\log n}\right),
\]
yielding a near-linear improvement in the number of APSP calls.

\paragraph{Concrete bounds assuming $k = O(n)$.}
We summarize the total time complexity for the two natural input scenarios, assuming the output spanner has $k = O(n)$ necessary edges.

\begin{itemize}
    \item \textbf{Complete graph input} ($m = \binom{n}{2} = O(n^2)$).
    Each APSP costs $O(n(n^2 + n \log n)) = O(n^3)$:
    \[
        O(n^2 \log n + n \log n \cdot n^3) = O(n^4 \log n).
    \]

    \item \textbf{$t$-spanner input} ($m = O(n)$).
    Each APSP costs $O(n(n + n \log n)) = O(n^2 \log n)$:
    \[
        O(n \log n + n \log n \cdot n^2 \log n) = O(n^3 \log^2 n).
    \]
\end{itemize}

\noindent When the input is already a $t$-spanner with $O(n)$ edges, the total complexity improves from $O(n^4 \log n)$ to $O(n^3 \log^2 n)$.


\section{Related Work}

Geometric spanners have been extensively studied since the 1980s. We survey the main constructions, their properties, and their relationship to our work. For a comprehensive treatment, we refer the reader to the monograph by Narasimhan and Smid~\cite{narasimhan2007geometric} and the tutorial survey by Ahmed et al.~\cite{ahmed2020graph}.

\subsection{The Greedy Spanner}

The greedy (or path-greedy) spanner was introduced by Alth\"ofer, Das, Dobkin, Joseph, and Soares~\cite{althofer1993sparse}, based on earlier work by Das and Joseph~\cite{das1989which}. The algorithm processes all $\binom{n}{2}$ pairs of points in non-decreasing order of distance and adds an edge $(u,v)$ if and only if the current graph does not already contain a $t$-short path between $u$ and $v$. For Euclidean point sets in $\mathbb{R}^d$ with stretch parameter $t = 1 + \varepsilon$, the greedy spanner achieves $O(\varepsilon^{-(d-1)} \cdot n)$ edges and total weight $O(\varepsilon^{-d} \cdot \mathrm{wt}(\mathrm{MST}))$. Le and Solomon~\cite{le2019truly} proved these bounds are tight: any $(1+\varepsilon)$-spanner requires $\Omega(\varepsilon^{-(d-1)} \cdot n)$ edges and lightness $\Omega(\varepsilon^{-d})$. The greedy spanner also has bounded degree and always contains the minimum spanning tree.

Filtser and Solomon~\cite{filtser2020greedy} proved that the greedy spanner is \emph{existentially optimal}, meaning its worst-case performance over any graph family matches (up to constants) the worst-case of an optimal spanner for that family. However, Le et al.~\cite{le2024instance} recently showed that the greedy spanner is far from \emph{instance-optimal}: for any specific input, it may use significantly more edges and weight than necessary.

The construction time of the greedy spanner has also been studied. Bose, Carmi, Farshi, Maheshwari, and Smid~\cite{bose2010computing} achieved near-quadratic $O(n^2 \log n)$ time for the exact greedy spanner, which is near-optimal given the $\Omega(n^2)$ lower bound. Gudmundsson, Levcopoulos, and Narasimhan~\cite{gudmundsson2002fast} gave $O(n \log n)$-time algorithms for approximate greedy spanners.

Our Descending Greedy Filter can be viewed as a ``dual'' of the greedy spanner: while the greedy algorithm builds up the spanner by adding edges from shortest to longest, DGF starts with all edges and removes them from longest to shortest. Both constructions produce edge-minimal spanners, and when $t = n - 1$, DGF produces the MST---mirroring the relationship between Kruskal's algorithm (which adds edges in ascending order) and the reverse-delete algorithm (which removes edges in descending order).

\subsection{Weight Bounds and the Leapfrog Property}

A central result in the theory of geometric spanners is that the greedy spanner has total weight $O(\mathrm{wt}(\mathrm{MST}))$ for fixed $t$ and $d$. This was established through the \emph{$t$-leapfrog property}, introduced and analyzed by Das, Narasimhan, and Salowe (see~\cite{narasimhan2007geometric}, Chapter~14). The leapfrog property provides a geometric condition on edge sets that implies bounded total weight relative to the MST. Chandra, Das, Narasimhan, and Soares~\cite{chandra1992new} proved related sparseness and weight results for graph spanners.

Further weight and lightness bounds were obtained by Arya, Das, Mount, Salowe, and Smid~\cite{arya1995euclidean}, who constructed $(1+\varepsilon)$-spanners with bounded degree, $O(\log n)$ diameter, and weight $O(\log^2 n \cdot \mathrm{wt}(\mathrm{MST}))$. Elkin and Solomon~\cite{elkin2013optimal} improved the weight to $O(\log n \cdot \mathrm{wt}(\mathrm{MST}))$, and Le and Solomon~\cite{le2023unified} later provided a unified framework for obtaining light spanners across many graph classes.

We conjecture that DGF also achieves weight $O(\mathrm{wt}(\mathrm{MST}))$, and more broadly that any $t$-minimal spanner has this property. Our proof that DGF produces the MST when $t = n-1$ provides initial evidence toward this conjecture.

\subsection{Edge-Minimality and Optimal Spanners}

The concept of \emph{edge-minimality} (or $t$-minimality) that we study---where removing any single edge from the spanner violates the $t$-spanner condition for at least one pair---appears to be largely unexplored in the literature as a first-class structural property. Prior work has focused instead on \emph{minimum} spanners (globally minimizing edge count or weight), which is a strictly stronger requirement. Cai~\cite{cai1994np} proved that the minimum $t$-spanner problem is NP-hard for fixed $t \geq 2$ on general graphs. Carmi and Chaitman-Yerushalmi~\cite{carmi2013minimum} extended this to the Euclidean setting, proving that finding a minimum-weight Euclidean $t$-spanner is NP-hard for every constant $t > 1$.

Exact computation of minimum-weight spanners has been studied by Sigurd and Zachariasen~\cite{sigurd2004construction}, who formulated the problem as an ILP and solved instances with up to 64 vertices, and more recently by B\"okler, Chimani, Jasper, and Wagner~\cite{bokler2024exact}, who scaled the approach to over 16,000 vertices using column generation.

While a minimum-weight (or minimum-edge) spanner is necessarily edge-minimal, the converse does not hold. Our DGF algorithm guarantees edge-minimality---a weaker but efficiently achievable property---which we believe still implies strong structural guarantees such as bounded weight.

\subsection{Variants and Generalizations of the Greedy Spanner}

Several variants of the greedy spanner have been proposed. The \emph{gap-greedy} algorithm, described in~\cite{narasimhan2007geometric} and further developed by Arya and Smid~\cite{arya1997efficient}, replaces exact shortest-path checks with a geometric gap criterion, leading to faster construction time. Abu-Affash, Bar-On, and Carmi~\cite{abuaffash2021delta} introduced the \emph{$\delta$-greedy} $t$-spanner, which interpolates between the path-greedy ($\delta = 0$) and gap-greedy algorithms by controlling the number of shortest-path queries, achieving the same theoretical guarantees with improved practical construction time. Bose, Carmi, Maheshwari and others~\cite{bose2023path} studied structural properties of the path-greedy spanner on convex point sets.

A common practical optimization is to first compute a $(t/t')$-spanner (e.g., using a $\Theta$-graph or WSPD), and then apply the greedy algorithm with parameter~$t'$ on this sparse subgraph. While this reduces the number of candidate edges, as we discussed in Section~3, it does not preserve edge-minimality in general.

\subsection{Experimental Comparisons}

Farshi and Gudmundsson~\cite{farshi2009experimental} conducted the most comprehensive experimental study of geometric $t$-spanners, comparing the greedy spanner, $\Theta$-graphs, ordered $\Theta$-graphs, WSPD spanners, skip-list spanners, and sink spanners across multiple quality measures including edge count, total weight, maximum degree, spanner diameter, and edge crossings. They found that the greedy spanner consistently dominates on quality measures but is the slowest to construct. Chimani and Stutzenstein~\cite{chimani2022spanner} provided the first experimental comparison of polynomial-time spanner approximation algorithms for general weighted graphs, finding that the simplest greedy approach (ADDJS) produces the sparsest and lightest spanners in practice. Ghosh, Maheshwari, and others~\cite{ghosh2023bounded} compared all eleven known bounded-degree plane geometric spanner algorithms experimentally.

Our experimental results contribute to this body of work by demonstrating that, on random point sets, the DGF construction consistently produces spanners with fewer edges and lower total weight than the greedy spanner---the first such result to our knowledge.

\subsection{Connection to the Reverse-Delete Algorithm for MST}

The classical reverse-delete algorithm for minimum spanning trees~\cite{kruskal1956shortest} processes edges in non-increasing weight order, removing each edge if the graph remains connected. It is dual to Kruskal's algorithm, and both are correct because spanning trees form a matroid. Our DGF algorithm generalizes the reverse-delete approach from connectivity (MST) to the $t$-spanner condition: it processes edges from longest to shortest, removing each edge whose removal does not violate the $t$-spanner property. When $t = n - 1$, the $t$-spanner condition degenerates to connectivity, and DGF reduces exactly to the reverse-delete MST algorithm, as we prove in Section~5. However, unlike spanning trees, the collection of $t$-spanners does not form a matroid, so the forward greedy and reverse-delete approaches generally produce different results. Understanding the structural differences between the two constructions is an interesting direction for future work.


\bibliographystyle{plain}
\begin{thebibliography}{99}

\bibitem{abuaffash2021delta}
A.~K. Abu-Affash, G.~Bar-On, and P.~Carmi.
\newblock $\delta$-Greedy $t$-spanner.
\newblock {\em Computational Geometry: Theory and Applications}, 99:101807, 2021.

\bibitem{ahmed2020graph}
R.~Ahmed, G.~Bodwin, F.~Sahneh, K.~Kobourov, and R.~Spence.
\newblock Graph spanners: A tutorial review.
\newblock {\em Computer Science Review}, 37:100253, 2020.

\bibitem{althofer1993sparse}
I.~Alth\"ofer, G.~Das, D.~Dobkin, D.~Joseph, and J.~Soares.
\newblock On sparse spanners of weighted graphs.
\newblock {\em Discrete \& Computational Geometry}, 9(1):81--100, 1993.

\bibitem{arya1997efficient}
S.~Arya and M.~Smid.
\newblock Efficient construction of a bounded-degree spanner with low weight.
\newblock {\em Algorithmica}, 17:33--54, 1997.

\bibitem{arya1995euclidean}
S.~Arya, G.~Das, D.~M. Mount, J.~S. Salowe, and M.~Smid.
\newblock Euclidean spanners: Short, thin, and lanky.
\newblock In {\em Proc.\ 27th ACM Symposium on Theory of Computing (STOC)}, pages 489--498, 1995.

\bibitem{bokler2024exact}
F.~B\"okler, M.~Chimani, H.~Jasper, and M.~H. Wagner.
\newblock Exact minimum weight spanners via column generation.
\newblock In {\em Proc.\ 32nd European Symposium on Algorithms (ESA)}, LIPIcs 308:30:1--30:17, 2024.

\bibitem{bose2010computing}
P.~Bose, P.~Carmi, M.~Farshi, A.~Maheshwari, and M.~Smid.
\newblock Computing the greedy spanner in near-quadratic time.
\newblock {\em Algorithmica}, 58(3):711--729, 2010.

\bibitem{bose2023path}
P.~Bose, P.~Carmi, A.~Maheshwari, and M.~Smid.
\newblock On path-greedy geometric spanners.
\newblock {\em Computational Geometry: Theory and Applications}, 2023.

\bibitem{cai1994np}
L.~Cai.
\newblock NP-completeness of minimum spanner problems.
\newblock {\em Discrete Applied Mathematics}, 48(2):187--194, 1994.

\bibitem{carmi2013minimum}
P.~Carmi and L.~Chaitman-Yerushalmi.
\newblock Minimum weight {E}uclidean $t$-spanner is {NP}-hard.
\newblock {\em Journal of Discrete Algorithms}, 22:30--42, 2013.

\bibitem{chandra1992new}
B.~Chandra, G.~Das, G.~Narasimhan, and J.~Soares.
\newblock New sparseness results on graph spanners.
\newblock {\em International Journal of Computational Geometry \& Applications}, 5:125--144, 1995.
\newblock Preliminary version in {\em Proc.\ 8th Annual ACM Symposium on Computational Geometry}, 1992.

\bibitem{chimani2022spanner}
M.~Chimani and F.~Stutzenstein.
\newblock Spanner approximations in practice.
\newblock In {\em Proc.\ 30th European Symposium on Algorithms (ESA)}, LIPIcs 244:37:1--37:15, 2022.

\bibitem{das1989which}
G.~Das and D.~Joseph.
\newblock Which triangulations approximate the complete graph?
\newblock In {\em Optimal Algorithms}, LNCS 401, pages 168--192. Springer, 1989.

\bibitem{elkin2013optimal}
M.~Elkin and S.~Solomon.
\newblock Optimal {E}uclidean spanners: Really short, thin, and lanky.
\newblock In {\em Proc.\ 45th ACM Symposium on Theory of Computing (STOC)}, pages 645--654, 2013.

\bibitem{farshi2009experimental}
M.~Farshi and J.~Gudmundsson.
\newblock Experimental study of geometric $t$-spanners.
\newblock {\em Journal of Experimental Algorithmics}, 14, 2009.

\bibitem{filtser2020greedy}
A.~Filtser and S.~Solomon.
\newblock The greedy spanner is existentially optimal.
\newblock {\em SIAM Journal on Computing}, 49(2):429--447, 2020.

\bibitem{ghosh2023bounded}
A.~Ghosh, A.~Maheshwari, et~al.
\newblock Bounded-degree plane geometric spanners in practice.
\newblock {\em ACM Journal of Experimental Algorithmics}, 2023.

\bibitem{gudmundsson2002fast}
J.~Gudmundsson, C.~Levcopoulos, and G.~Narasimhan.
\newblock Fast greedy algorithms for constructing sparse geometric spanners.
\newblock {\em SIAM Journal on Computing}, 31(5), 2002.

\bibitem{kruskal1956shortest}
J.~B. Kruskal.
\newblock On the shortest spanning subtree of a graph and the traveling salesman problem.
\newblock {\em Proceedings of the American Mathematical Society}, 7(1):48--50, 1956.

\bibitem{le2024instance}
H.~Le, S.~Solomon, C.~Than, C.~D. T\'oth, and T.~Zhang.
\newblock Towards instance-optimal {E}uclidean spanners.
\newblock In {\em Proc.\ 65th IEEE Symposium on Foundations of Computer Science (FOCS)}, 2024.

\bibitem{le2019truly}
H.~Le and S.~Solomon.
\newblock Truly optimal {E}uclidean spanners.
\newblock In {\em Proc.\ 60th IEEE Symposium on Foundations of Computer Science (FOCS)}, pages 1078--1100, 2019.

\bibitem{le2023unified}
H.~Le and S.~Solomon.
\newblock A unified framework for light spanners.
\newblock In {\em Proc.\ 55th ACM Symposium on Theory of Computing (STOC)}, pages 295--308, 2023.

\bibitem{narasimhan2007geometric}
G.~Narasimhan and M.~Smid.
\newblock {\em Geometric Spanner Networks}.
\newblock Cambridge University Press, 2007.

\bibitem{sigurd2004construction}
M.~Sigurd and M.~Zachariasen.
\newblock Construction of minimum-weight spanners.
\newblock In {\em Proc.\ 12th European Symposium on Algorithms (ESA)}, LNCS 3221, pages 797--808, 2004.

\end{thebibliography}

\end{document}
